\documentclass[12pt,a4paper]{article}
\usepackage[utf8]{inputenc}
\usepackage[T1]{fontenc}
\usepackage[french]{babel}
\usepackage[margin=2.5cm]{geometry}
\usepackage{xcolor}
\usepackage{listings}
\usepackage{hyperref}
\usepackage{booktabs}
\usepackage{tabularx}
\usepackage{longtable}
\usepackage{array}
\usepackage{fancyhdr}
\usepackage{titlesec}
\usepackage{graphicx}
\usepackage{mdframed}
\usepackage{enumitem}
\usepackage{amsmath}
\usepackage{amssymb}

% ---- Couleurs ----
\definecolor{primary}{RGB}{26, 86, 219}
\definecolor{secondary}{RGB}{30, 64, 175}
\definecolor{success}{RGB}{5, 150, 105}
\definecolor{danger}{RGB}{220, 38, 38}
\definecolor{warning}{RGB}{217, 119, 6}
\definecolor{codebg}{RGB}{245, 247, 250}
\definecolor{codeborder}{RGB}{203, 213, 225}
\definecolor{codecomment}{RGB}{107, 114, 128}
\definecolor{codestring}{RGB}{5, 150, 105}
\definecolor{codekeyword}{RGB}{99, 102, 241}
\definecolor{lightblue}{RGB}{239, 246, 255}
\definecolor{lightgreen}{RGB}{236, 253, 245}
\definecolor{lightred}{RGB}{254, 242, 242}

% ---- Style des listings ----
\lstdefinestyle{python}{
  language=Python,
  backgroundcolor=\color{codebg},
  basicstyle=\ttfamily\footnotesize,
  keywordstyle=\color{codekeyword}\bfseries,
  commentstyle=\color{codecomment}\itshape,
  stringstyle=\color{codestring},
  numberstyle=\tiny\color{codecomment},
  numbers=left,
  stepnumber=1,
  numbersep=8pt,
  frame=single,
  rulecolor=\color{codeborder},
  breaklines=true,
  breakatwhitespace=true,
  tabsize=4,
  showstringspaces=false,
  xleftmargin=15pt,
  xrightmargin=5pt,
  captionpos=b,
}

\lstdefinestyle{json}{
  language=,
  backgroundcolor=\color{codebg},
  basicstyle=\ttfamily\footnotesize,
  keywordstyle=\color{codekeyword},
  stringstyle=\color{codestring},
  frame=single,
  rulecolor=\color{codeborder},
  breaklines=true,
  tabsize=2,
  showstringspaces=false,
  xleftmargin=15pt,
  captionpos=b,
}

\lstdefinestyle{bash}{
  language=bash,
  backgroundcolor=\color{codebg},
  basicstyle=\ttfamily\footnotesize\color{white},
  backgroundcolor=\color[RGB]{30,41,59},
  keywordstyle=\color[RGB]{129,230,217},
  stringstyle=\color[RGB]{167,243,208},
  commentstyle=\color[RGB]{107,114,128}\itshape,
  frame=single,
  rulecolor=\color[RGB]{51,65,85},
  breaklines=true,
  tabsize=2,
  showstringspaces=false,
  xleftmargin=15pt,
  captionpos=b,
}

% ---- En-tête et pied de page ----
\pagestyle{fancy}
\fancyhf{}
\fancyhead[L]{\small\textcolor{primary}{\textbf{Devoir 304 — API de transaction bancaire}}}
\fancyhead[R]{\small\textcolor{codecomment}{KOA ESSAMA PAULIN BRICE}}
\fancyfoot[C]{\small\textcolor{codecomment}{\thepage}}
\renewcommand{\headrulewidth}{0.5pt}
\renewcommand{\footrulewidth}{0.5pt}

% ---- Titres ----
\titleformat{\section}{\Large\bfseries\color{primary}}{}{0em}{\thesection\quad}[\titlerule]
\titleformat{\subsection}{\large\bfseries\color{secondary}}{}{0em}{\thesubsection\quad}
\titleformat{\subsubsection}{\normalsize\bfseries\color{secondary}}{}{0em}{\thesubsubsection\quad}

% ---- Boîtes colorées ----
\newmdenv[
  backgroundcolor=lightblue,
  linecolor=primary,
  linewidth=2pt,
  leftline=true, rightline=false, topline=false, bottomline=false,
  innerleftmargin=10pt,
  innerrightmargin=10pt,
  innertopmargin=8pt,
  innerbottommargin=8pt,
]{infobox}

\newmdenv[
  backgroundcolor=lightgreen,
  linecolor=success,
  linewidth=2pt,
  leftline=true, rightline=false, topline=false, bottomline=false,
  innerleftmargin=10pt,
  innerrightmargin=10pt,
  innertopmargin=8pt,
  innerbottommargin=8pt,
]{successbox}

\newmdenv[
  backgroundcolor=lightred,
  linecolor=danger,
  linewidth=2pt,
  leftline=true, rightline=false, topline=false, bottomline=false,
  innerleftmargin=10pt,
  innerrightmargin=10pt,
  innertopmargin=8pt,
  innerbottommargin=8pt,
]{dangerbox}

\hypersetup{
  colorlinks=true,
  linkcolor=primary,
  urlcolor=primary,
  pdftitle={Rapport d'implementation - API Bancaire},
  pdfauthor={KOA ESSAMA PAULIN BRICE},
}

\setlength{\parskip}{6pt}
\setlength{\parindent}{0pt}

% =========================================================
\begin{document}
% =========================================================

% ---- Page de garde ----
\begin{titlepage}
\centering
\vspace*{1cm}

{\Large\bfseries UNIVERSITÉ DE YAOUNDÉ I\\[0.3cm]}
{\large Faculté des Sciences\\[0.1cm]}
{\large Département d'Informatique\\[2cm]}

\noindent\rule{\linewidth}{2pt}\\[0.5cm]
{\Huge\bfseries\color{primary} Rapport d'Implémentation}\\[0.3cm]
{\LARGE\bfseries API de Transaction Bancaire}\\[0.3cm]
{\large Devoir 304 — Conception d'API REST}\\[0.3cm]
\noindent\rule{\linewidth}{2pt}\\[2cm]

\begin{tabular}{ll}
\textbf{Étudiant :} & KOA ESSAMA PAULIN BRICE \\[0.3cm]
\textbf{Matricule :} & 25I2255 \\[0.3cm]
\textbf{Niveau :} & Licence 3 ICT \\[0.3cm]
\textbf{Établissement :} & Université de Yaoundé I \\[0.3cm]
\textbf{Dépôt GitHub :} & \url{https://github.com/eliote-geeks/Devoir_304} \\[0.3cm]
\textbf{Date :} & Avril 2026 \\
\end{tabular}

\vfill
{\large\textit{«~Conception et implémentation d'une API REST pour la gestion\\de comptes bancaires, de dépôts et de retraits.~»}}

\end{titlepage}

% ---- Table des matières ----
\tableofcontents
\newpage

% =========================================================
\section{Introduction et contexte}
% =========================================================

\subsection{Contexte du projet}

Ce rapport décrit de façon complète et détaillée l'implémentation d'une \textbf{API REST de transaction bancaire} réalisée dans le cadre du Devoir 304 de Licence 3 ICT.

L'objectif est de concevoir un système permettant à des clients de gérer leurs comptes bancaires numériquement : créer un compte, consulter son solde, effectuer des dépôts et des retraits, et visualiser l'historique de leurs opérations.

\subsection{Problème résolu}

Les établissements financiers ont besoin de solutions numériques pour gérer les opérations de base. Cette API constitue le back-end (côté serveur) d'un tel système, pouvant être consommée par une application web, mobile ou desktop.

\subsection{Choix technologiques}

\begin{center}
\begin{tabular}{lll}
\toprule
\textbf{Composant} & \textbf{Technologie} & \textbf{Raison du choix} \\
\midrule
Langage & Python 3.13 & Lisible, productif, écosystème riche \\
Framework & FastAPI & Moderne, rapide, génère Swagger automatiquement \\
Serveur & Uvicorn & Serveur ASGI haute performance \\
Persistance & Fichier JSON local & Simple, portable, sans dépendance externe \\
Documentation & Swagger UI (OpenAPI 3.0) & Standard industriel, interactif \\
Validation & Pydantic v2 & Validation automatique des données entrantes \\
\bottomrule
\end{tabular}
\end{center}

% =========================================================
\section{Architecture du projet}
% =========================================================

\subsection{Structure des fichiers}

\begin{lstlisting}[style=bash, caption={Arborescence du projet}]
Devoir_304/
|
|-- app/                      # Package principal de l'API
|   |-- __init__.py           # Marque app/ comme package Python
|   |-- main.py               # Serveur FastAPI : routes et logique HTTP
|   |-- models.py             # Schemas de validation (Pydantic)
|   |-- storage.py            # Persistance : lecture/ecriture JSON
|   |-- static/               # Fichiers statiques (Swagger UI)
|       |-- swagger-ui/
|           |-- swagger-ui-bundle.js
|           |-- swagger-ui.css
|           |-- favicon.svg
|
|-- data/
|   |-- accounts.json         # Base de donnees locale (fichier JSON)
|
|-- requirements.txt          # Dependances Python
|-- render.yaml               # Configuration deploiement Render
|-- CAHIER_DE_CHARGES.md      # Cahier des charges complet
\end{lstlisting}

\subsection{Architecture en couches}

L'API est organisée en \textbf{3 couches} distinctes, selon le principe de séparation des responsabilités :

\begin{center}
\begin{tabular}{clp{8cm}}
\toprule
\textbf{Couche} & \textbf{Fichier} & \textbf{Rôle} \\
\midrule
\textbf{Présentation} & \texttt{main.py} & Définit les routes HTTP, reçoit les requêtes, retourne les réponses JSON \\
\textbf{Modèle} & \texttt{models.py} & Définit les structures de données et règles de validation \\
\textbf{Données} & \texttt{storage.py} & Gère la lecture et l'écriture dans le fichier \texttt{accounts.json} \\
\bottomrule
\end{tabular}
\end{center}

\begin{infobox}
\textbf{Pourquoi cette séparation ?} Elle permet de modifier chaque couche indépendamment. Par exemple, on pourrait remplacer le fichier JSON par une vraie base de données (PostgreSQL) en modifiant uniquement \texttt{storage.py}, sans toucher aux routes ni aux modèles.
\end{infobox}

% =========================================================
\section{Spécifications fonctionnelles}
% =========================================================

\subsection{Acteurs du système}

\textbf{Client bancaire} : l'utilisateur final qui interagit avec le système. Il peut :
\begin{itemize}[noitemsep]
  \item ouvrir un compte bancaire
  \item consulter les informations de son compte
  \item effectuer des dépôts
  \item effectuer des retraits
  \item consulter son historique de transactions
\end{itemize}

\textbf{Administrateur système} : supervise le système. Il peut :
\begin{itemize}[noitemsep]
  \item consulter la liste complète des comptes
  \item suivre les transactions
\end{itemize}

\subsection{Les 5 fonctions du système}

\begin{center}
\begin{tabular}{clll}
\toprule
\textbf{\#} & \textbf{Fonction} & \textbf{Méthode HTTP} & \textbf{URL} \\
\midrule
F1 & Créer un compte & POST & \texttt{/accounts} \\
F2 & Lister les comptes & GET & \texttt{/accounts} \\
F3 & Consulter un compte & GET & \texttt{/accounts/\{id\}} \\
F4 & Effectuer un dépôt & POST & \texttt{/accounts/\{id\}/deposit} \\
F5 & Effectuer un retrait & POST & \texttt{/accounts/\{id\}/withdraw} \\
F+ & Historique des transactions & GET & \texttt{/accounts/\{id\}/transactions} \\
\bottomrule
\end{tabular}
\end{center}

\subsection{Règles de gestion}

\begin{enumerate}[noitemsep]
  \item Un compte possède un identifiant unique (UUID généré automatiquement)
  \item Un compte possède un numéro de compte unique au format \texttt{ACC-AAAAMMJJ-XXXXXXXX}
  \item Le solde initial doit être supérieur ou égal à zéro
  \item Un dépôt doit obligatoirement avoir un montant strictement positif
  \item Un retrait doit obligatoirement avoir un montant strictement positif
  \item Un retrait est refusé si le solde du compte est insuffisant
  \item Chaque opération modifie immédiatement et définitivement le solde
  \item Chaque opération est enregistrée dans l'historique avec sa date et le solde après
\end{enumerate}

% =========================================================
\section{Spécifications non fonctionnelles}
% =========================================================

\begin{center}
\begin{tabular}{lp{4.5cm}p{6cm}}
\toprule
\textbf{Critère} & \textbf{Exigence} & \textbf{Comment c'est respecté} \\
\midrule
Performance & Réponses rapides & FastAPI + Uvicorn (ASGI asynchrone), pas de base de données distante \\
Fiabilité & Données persistées & Fichier \texttt{accounts.json} protégé par un verrou \texttt{threading.Lock} \\
Sécurité & Montants invalides bloqués & Pydantic valide \texttt{gt=0} et \texttt{ge=0} avant tout traitement \\
Maintenabilité & Code modulaire & 3 fichiers avec rôles clairs (\texttt{main}, \texttt{models}, \texttt{storage}) \\
Évolutivité & Extensible & Architecture prête pour authentification JWT, virements, base SQL \\
Disponibilité & API accessible & Fonctionne tant que le serveur Uvicorn est actif \\
\bottomrule
\end{tabular}
\end{center}

% =========================================================
\section{Implémentation détaillée — fichier par fichier}
% =========================================================

\subsection{\texttt{models.py} — Définition des structures de données}

Ce fichier définit tous les \textbf{schémas de données} utilisés par l'API. Chaque classe hérite de \texttt{BaseModel} de Pydantic, ce qui permet la validation automatique des données entrantes.

\subsubsection{AccountCreate — Données pour créer un compte}

\begin{lstlisting}[style=python, caption={models.py — AccountCreate}]
class AccountCreate(BaseModel):
    full_name: str = Field(
        ...,            # champ obligatoire (les ... signifient "requis")
        min_length=3,   # au moins 3 caracteres
        max_length=120,
        description="Nom complet du titulaire du compte.",
        examples=["KOA ESSAMA PAULIN BRICE"],
    )
    phone_number: str = Field(
        ...,            # champ obligatoire
        min_length=6,
        max_length=20,
        description="Numero de telephone du client.",
        examples=["699001122"],
    )
    email: str | None = Field(
        default=None,   # optionnel : si absent, vaut None
        max_length=120,
        description="Adresse email du client.",
        examples=["paulin@example.com"],
    )
    initial_balance: float = Field(
        default=0,      # optionnel : defaut = 0
        ge=0,           # ge = "greater or equal" : interdit les negatifs
        description="Solde initial. Doit etre superieur ou egal a zero.",
        examples=[10000],
    )
\end{lstlisting}

\begin{infobox}
\textbf{Comment Pydantic fonctionne :} Quand une requête POST arrive avec un corps JSON, Pydantic lit automatiquement les champs, vérifie les contraintes (\texttt{min\_length}, \texttt{ge}, etc.) et retourne une erreur HTTP 422 si les données sont invalides — sans que le développeur écrive une seule ligne de validation.
\end{infobox}

\subsubsection{TransactionRequest — Données pour un dépôt ou retrait}

\begin{lstlisting}[style=python, caption={models.py — TransactionRequest}]
class TransactionRequest(BaseModel):
    amount: float = Field(
        ...,    # obligatoire
        gt=0,   # gt = "greater than" : strictement positif
        description="Montant de l'operation. Doit etre strictement positif.",
        examples=[5000],
    )
    description: str | None = Field(
        default=None,   # optionnel : libelle de l'operation
        max_length=150,
        description="Libelle ou motif de la transaction.",
        examples=["Depot agence"],
    )
\end{lstlisting}

\subsubsection{Autres modèles — Réponses de l'API}

\begin{lstlisting}[style=python, caption={models.py — Modeles de reponse}]
class AccountResponse(BaseModel):
    id: str                   # UUID du compte
    account_number: str       # ex: ACC-20260422-A1B2C3D4
    full_name: str
    phone_number: str
    email: str | None
    balance: float            # solde courant
    created_at: str           # date ISO 8601

class TransactionResponse(BaseModel):
    transaction_id: str
    transaction_type: Literal["deposit", "withdraw"]
    amount: float
    description: str | None
    balance_after: float      # solde apres l'operation
    created_at: str

class AccountDetails(AccountResponse):
    # Herite de AccountResponse et ajoute les transactions
    transactions: list[TransactionResponse] = Field(default_factory=list)

class ErrorResponse(BaseModel):
    detail: str               # Message d'erreur retourne par l'API
\end{lstlisting}

% -------------------------------------------------------
\subsection{\texttt{storage.py} — Persistance des données}
% -------------------------------------------------------

Ce fichier gère toute la \textbf{lecture et écriture des données} dans le fichier JSON local.

\subsubsection{Configuration et utilitaires}

\begin{lstlisting}[style=python, caption={storage.py — Configuration}]
import json
from pathlib import Path
from threading import Lock    # protege les ecritures concurrentes
from uuid import uuid4
from datetime import datetime, UTC

BASE_DIR = Path(__file__).resolve().parent.parent   # racine du projet
DATA_FILE = BASE_DIR / "data" / "accounts.json"    # chemin du fichier JSON
FILE_LOCK = Lock()    # verrou pour empecher 2 ecritures simultanees

def _now() -> str:
    # Retourne la date/heure actuelle au format ISO 8601
    # Ex: "2026-04-22T14:30:00Z"
    return datetime.now(UTC).isoformat().replace("+00:00", "Z")

def _generate_account_number() -> str:
    # Genere un numero de compte unique
    # Ex: "ACC-20260422-A1B2C3D4"
    date_part = datetime.now(UTC).strftime("%Y%m%d")
    random_part = uuid4().hex[:8].upper()   # 8 caracteres aleatoires
    return f"ACC-{date_part}-{random_part}"
\end{lstlisting}

\subsubsection{Initialisation et lecture/écriture}

\begin{lstlisting}[style=python, caption={storage.py — Lecture et ecriture}]
def initialize_storage() -> None:
    # Cree le dossier data/ et le fichier accounts.json s'ils n'existent pas
    DATA_FILE.parent.mkdir(parents=True, exist_ok=True)
    if not DATA_FILE.exists():
        DATA_FILE.write_text(
            json.dumps({"accounts": []}, indent=2),
            encoding="utf-8"
        )

def _read_data() -> dict:
    # Lit et desserialise le fichier JSON
    initialize_storage()    # s'assure que le fichier existe
    return json.loads(DATA_FILE.read_text(encoding="utf-8"))

def _write_data(data: dict) -> None:
    # Serialise et ecrit les donnees dans le fichier JSON
    DATA_FILE.write_text(json.dumps(data, indent=2), encoding="utf-8")
\end{lstlisting}

\begin{infobox}
\textbf{Structure du fichier \texttt{accounts.json} :}\\
Le fichier contient un objet JSON avec une clé \texttt{"accounts"} dont la valeur est un tableau de tous les comptes. Chaque compte contient ses transactions intégrées directement.
\end{infobox}

\begin{lstlisting}[style=json, caption={Exemple de contenu de accounts.json}]
{
  "accounts": [
    {
      "id": "a3f1b2c4-...",
      "account_number": "ACC-20260422-A1B2C3D4",
      "full_name": "KOA ESSAMA PAULIN BRICE",
      "phone_number": "699001122",
      "email": "paulin@example.com",
      "balance": 8000.0,
      "created_at": "2026-04-22T10:00:00Z",
      "transactions": [
        {
          "transaction_id": "f1e2d3c4-...",
          "transaction_type": "deposit",
          "amount": 10000.0,
          "description": "Solde initial",
          "balance_after": 10000.0,
          "created_at": "2026-04-22T10:00:00Z"
        },
        {
          "transaction_id": "b5a6c7d8-...",
          "transaction_type": "withdraw",
          "amount": 2000.0,
          "description": "Retrait DAB",
          "balance_after": 8000.0,
          "created_at": "2026-04-22T11:30:00Z"
        }
      ]
    }
  ]
}
\end{lstlisting}

\subsubsection{Création d'un compte}

\begin{lstlisting}[style=python, caption={storage.py — create\_account()}]
def create_account(payload: dict) -> dict:
    with FILE_LOCK:     # verrou : une seule operation d'ecriture a la fois
        data = _read_data()

        account = {
            "id": str(uuid4()),                      # UUID unique
            "account_number": _generate_account_number(),
            "full_name": payload["full_name"],
            "phone_number": payload["phone_number"],
            "email": payload.get("email"),           # None si absent
            "balance": round(float(payload.get("initial_balance", 0)), 2),
            "created_at": _now(),
            "transactions": [],                      # liste vide au depart
        }

        # Si un solde initial est fourni, on cree une premiere transaction
        if account["balance"] > 0:
            account["transactions"].append({
                "transaction_id": str(uuid4()),
                "transaction_type": "deposit",
                "amount": account["balance"],
                "description": "Solde initial",
                "balance_after": account["balance"],
                "created_at": _now(),
            })

        data["accounts"].append(account)    # ajout dans la liste
        _write_data(data)                   # sauvegarde dans le fichier
        return account
\end{lstlisting}

\subsubsection{Application d'une transaction (dépôt ou retrait)}

\begin{lstlisting}[style=python, caption={storage.py — apply\_transaction()}]
def apply_transaction(
    account_id: str,
    transaction_type: str,   # "deposit" ou "withdraw"
    amount: float,
    description: str | None
) -> dict | None:

    with FILE_LOCK:     # verrou pour l'ecriture atomique
        data = _read_data()

        for account in data["accounts"]:
            if account["id"] != account_id:
                continue    # passe au compte suivant si ce n'est pas le bon

            current_balance = float(account["balance"])
            amount = round(float(amount), 2)

            # Verification du solde avant un retrait
            if transaction_type == "withdraw" and amount > current_balance:
                raise ValueError("Solde insuffisant pour effectuer ce retrait.")

            # Calcul du nouveau solde
            if transaction_type == "deposit":
                new_balance = round(current_balance + amount, 2)
            else:
                new_balance = round(current_balance - amount, 2)

            # Creation de la transaction
            transaction = {
                "transaction_id": str(uuid4()),
                "transaction_type": transaction_type,
                "amount": amount,
                "description": description,
                "balance_after": new_balance,
                "created_at": _now(),
            }

            account["balance"] = new_balance            # mise a jour du solde
            account["transactions"].append(transaction) # ajout dans l'historique
            _write_data(data)                           # sauvegarde
            return account

        return None     # compte non trouve
\end{lstlisting}

% -------------------------------------------------------
\subsection{\texttt{main.py} — Les routes de l'API}
% -------------------------------------------------------

Ce fichier est le cœur de l'API. Il définit toutes les \textbf{routes HTTP} et orchestre les appels aux autres modules.

\subsubsection{Initialisation de FastAPI}

\begin{lstlisting}[style=python, caption={main.py — Initialisation de l'application}]
from fastapi import FastAPI, HTTPException, status
from fastapi.openapi.docs import get_swagger_ui_html
from fastapi.staticfiles import StaticFiles

app = FastAPI(
    title="API de transaction bancaire",
    description="...",   # description longue visible dans Swagger
    version="1.0.0",
    docs_url=None,       # on desactive la route /docs par defaut
    redoc_url=None,      # pour utiliser notre propre Swagger UI local
    openapi_tags=[
        {"name": "General", "description": "Informations generales."},
        {"name": "Accounts", "description": "Gestion des comptes."},
        {"name": "Transactions", "description": "Depots et retraits."},
    ],
)

# Sert les fichiers du dossier app/static/ a l'URL /static/
app.mount("/static", StaticFiles(directory="app/static"), name="static")

@app.on_event("startup")
def startup() -> None:
    initialize_storage()    # cree le fichier JSON au demarrage si besoin
\end{lstlisting}

\subsubsection{Fonction 1 — Créer un compte}

\begin{lstlisting}[style=python, caption={main.py — POST /accounts}]
@app.post(
    "/accounts",
    response_model=AccountResponse,
    status_code=status.HTTP_201_CREATED,   # HTTP 201 Created (pas 200)
    tags=["Accounts"],
    summary="Creer un compte bancaire",
)
def create_bank_account(payload: AccountCreate) -> dict:
    # FastAPI injecte automatiquement payload apres validation Pydantic
    account = create_account(payload.model_dump())   # cree et persiste
    # On retourne le compte sans la liste de transactions
    return {key: value for key, value in account.items() if key != "transactions"}
\end{lstlisting}

\textbf{Ce qui se passe étape par étape :}
\begin{enumerate}[noitemsep]
  \item Le client envoie un POST avec un corps JSON
  \item FastAPI lit le corps et le passe à Pydantic (\texttt{AccountCreate})
  \item Pydantic vérifie les contraintes (champs obligatoires, longueurs, \texttt{ge=0})
  \item Si invalide → FastAPI retourne automatiquement HTTP 422 avec le détail
  \item Si valide → \texttt{create\_account()} est appelé, le compte est créé et sauvegardé
  \item La réponse HTTP 201 est retournée avec les données du nouveau compte
\end{enumerate}

\subsubsection{Fonction 2 — Lister les comptes}

\begin{lstlisting}[style=python, caption={main.py — GET /accounts}]
@app.get(
    "/accounts",
    response_model=list[AccountResponse],   # liste de comptes en reponse
    tags=["Accounts"],
    summary="Lister tous les comptes",
)
def get_accounts() -> list[dict]:
    accounts = list_accounts()    # lit tous les comptes depuis le fichier JSON
    # Retourne chaque compte sans ses transactions (pour alleges la reponse)
    return [
        {key: value for key, value in account.items() if key != "transactions"}
        for account in accounts
    ]
\end{lstlisting}

\subsubsection{Fonction 3 — Consulter un compte}

\begin{lstlisting}[style=python, caption={main.py — GET /accounts/\{account\_id\}}]
@app.get(
    "/accounts/{account_id}",   # {account_id} est un parametre de chemin
    response_model=AccountDetails,  # inclut les transactions
    tags=["Accounts"],
    summary="Consulter un compte",
    responses={404: {"model": ErrorResponse}},
)
def get_account_by_id(account_id: str) -> dict:
    account = get_account(account_id)   # cherche dans le fichier JSON
    if not account:
        # Lance une exception HTTP 404 si le compte n'existe pas
        raise HTTPException(status_code=404, detail="Compte introuvable.")
    return account   # retourne le compte avec ses transactions
\end{lstlisting}

\subsubsection{Fonction 4 — Effectuer un dépôt}

\begin{lstlisting}[style=python, caption={main.py — POST /accounts/\{account\_id\}/deposit}]
@app.post(
    "/accounts/{account_id}/deposit",
    response_model=AccountDetails,
    tags=["Transactions"],
    summary="Effectuer un depot",
    responses={
        404: {"model": ErrorResponse, "description": "Compte introuvable."},
    },
)
def deposit_money(account_id: str, payload: TransactionRequest) -> dict:
    # Appelle apply_transaction avec type "deposit"
    updated_account = apply_transaction(
        account_id=account_id,
        transaction_type="deposit",
        amount=payload.amount,
        description=payload.description,
    )
    if not updated_account:
        # apply_transaction retourne None si le compte n'existe pas
        raise HTTPException(status_code=404, detail="Compte introuvable.")
    return updated_account   # retourne le compte mis a jour
\end{lstlisting}

\subsubsection{Fonction 5 — Effectuer un retrait}

\begin{lstlisting}[style=python, caption={main.py — POST /accounts/\{account\_id\}/withdraw}]
@app.post(
    "/accounts/{account_id}/withdraw",
    response_model=AccountDetails,
    tags=["Transactions"],
    summary="Effectuer un retrait",
    responses={
        400: {"model": ErrorResponse, "description": "Solde insuffisant."},
        404: {"model": ErrorResponse, "description": "Compte introuvable."},
    },
)
def withdraw_money(account_id: str, payload: TransactionRequest) -> dict:
    try:
        updated_account = apply_transaction(
            account_id=account_id,
            transaction_type="withdraw",
            amount=payload.amount,
            description=payload.description,
        )
    except ValueError as error:
        # apply_transaction leve ValueError si solde insuffisant
        # On la convertit en erreur HTTP 400
        raise HTTPException(status_code=400, detail=str(error)) from error

    if not updated_account:
        raise HTTPException(status_code=404, detail="Compte introuvable.")
    return updated_account
\end{lstlisting}

% =========================================================
\section{Cas de test}
% =========================================================

\subsection{Tableau complet des 23 cas de test}

\begin{longtable}{clp{3.5cm}p{3cm}c}
\toprule
\textbf{ID} & \textbf{Fonction} & \textbf{Description} & \textbf{Résultat attendu} & \textbf{Code HTTP} \\
\midrule
\endhead
CT-01 & Créer compte & Tous les champs fournis & Compte créé avec ID et numéro & 201 \\
CT-02 & Créer compte & Sans email ni solde initial & Compte créé, balance = 0 & 201 \\
CT-03 & Créer compte & \texttt{full\_name} absent & Erreur de validation & 422 \\
CT-04 & Créer compte & \texttt{phone\_number} absent & Erreur de validation & 422 \\
CT-05 & Créer compte & Solde initial négatif & Erreur de validation & 422 \\
\midrule
CT-06 & Lister comptes & Aucun compte créé & Tableau vide \texttt{[]} & 200 \\
CT-07 & Lister comptes & Comptes existants & Liste des comptes & 200 \\
\midrule
CT-08 & Consulter compte & ID valide & Détails + transactions & 200 \\
CT-09 & Consulter compte & ID invalide & Compte introuvable & 404 \\
\midrule
CT-10 & Dépôt & Montant valide & Solde augmente & 200 \\
CT-11 & Dépôt & Avec description & Libellé visible & 200 \\
CT-12 & Dépôt & Montant = 0 & Erreur validation & 422 \\
CT-13 & Dépôt & Montant négatif & Erreur validation & 422 \\
CT-14 & Dépôt & Compte inexistant & Compte introuvable & 404 \\
\midrule
CT-15 & Retrait & Solde suffisant & Solde diminue & 200 \\
CT-16 & Retrait & Montant = solde & Solde = 0 & 200 \\
CT-17 & Retrait & Solde insuffisant & Solde insuffisant & 400 \\
CT-18 & Retrait & Montant = 0 & Erreur validation & 422 \\
CT-19 & Retrait & Montant négatif & Erreur validation & 422 \\
CT-20 & Retrait & Compte inexistant & Compte introuvable & 404 \\
\midrule
CT-21 & Transactions & Compte sans opération & Tableau vide & 200 \\
CT-22 & Transactions & Avec opérations & Liste chronologique & 200 \\
CT-23 & Transactions & Compte inexistant & Compte introuvable & 404 \\
\bottomrule
\end{longtable}

\subsection{Test manuel détaillé — Scénario complet}

Le test manuel se fait directement dans \textbf{Swagger UI} accessible à \texttt{http://localhost:8000/docs}.

\subsubsection{Étape 1 — Vérifier la liste vide (CT-06)}

\begin{lstlisting}[style=bash, caption={Appel GET /accounts}]
GET http://localhost:8000/accounts
\end{lstlisting}

\begin{successbox}
\textbf{Résultat attendu :} HTTP 200 avec le corps \texttt{[]}\\
Cela confirme que la persistance est initialisée et que le système fonctionne.
\end{successbox}

\subsubsection{Étape 2 — Créer un compte (CT-01)}

\begin{lstlisting}[style=bash, caption={Appel POST /accounts}]
POST http://localhost:8000/accounts
Content-Type: application/json

{
  "full_name": "KOA ESSAMA PAULIN BRICE",
  "phone_number": "699001122",
  "email": "paulin@example.com",
  "initial_balance": 10000
}
\end{lstlisting}

\begin{lstlisting}[style=json, caption={Reponse attendue — HTTP 201}]
{
  "id": "a3f1b2c4-d5e6-7890-abcd-ef1234567890",
  "account_number": "ACC-20260422-A1B2C3D4",
  "full_name": "KOA ESSAMA PAULIN BRICE",
  "phone_number": "699001122",
  "email": "paulin@example.com",
  "balance": 10000.0,
  "created_at": "2026-04-22T10:00:00Z"
}
\end{lstlisting}

\begin{infobox}
\textbf{À noter :} Copier l'\texttt{id} retourné pour l'utiliser dans les étapes suivantes.
\end{infobox}

\subsubsection{Étape 3 — Tester la validation (CT-03, CT-04, CT-05)}

\begin{lstlisting}[style=bash, caption={Test CT-05 : solde negatif}]
POST http://localhost:8000/accounts
Content-Type: application/json

{
  "full_name": "TEST",
  "phone_number": "600000001",
  "initial_balance": -500
}
\end{lstlisting}

\begin{dangerbox}
\textbf{Résultat attendu :} HTTP 422 avec un message d'erreur Pydantic indiquant que \texttt{initial\_balance} doit être $\geq 0$.
\end{dangerbox}

\subsubsection{Étape 4 — Effectuer un dépôt (CT-10)}

\begin{lstlisting}[style=bash, caption={Appel POST /accounts/\{id\}/deposit}]
POST http://localhost:8000/accounts/a3f1b2c4-.../deposit
Content-Type: application/json

{
  "amount": 5000,
  "description": "Depot agence"
}
\end{lstlisting}

\begin{lstlisting}[style=json, caption={Reponse attendue — HTTP 200, balance = 15000}]
{
  "id": "a3f1b2c4-...",
  "balance": 15000.0,
  "transactions": [
    {
      "transaction_type": "deposit",
      "amount": 10000.0,
      "description": "Solde initial",
      "balance_after": 10000.0
    },
    {
      "transaction_type": "deposit",
      "amount": 5000.0,
      "description": "Depot agence",
      "balance_after": 15000.0
    }
  ]
}
\end{lstlisting}

\subsubsection{Étape 5 — Effectuer un retrait valide (CT-15)}

\begin{lstlisting}[style=bash, caption={Appel POST /accounts/\{id\}/withdraw}]
POST http://localhost:8000/accounts/a3f1b2c4-.../withdraw
Content-Type: application/json

{
  "amount": 2000,
  "description": "Retrait DAB"
}
\end{lstlisting}

\begin{successbox}
\textbf{Résultat attendu :} HTTP 200, balance = 13000 (15000 - 2000).
\end{successbox}

\subsubsection{Étape 6 — Tester le solde insuffisant (CT-17)}

\begin{lstlisting}[style=bash, caption={Test CT-17 : retrait impossible}]
POST http://localhost:8000/accounts/a3f1b2c4-.../withdraw
Content-Type: application/json

{
  "amount": 50000,
  "description": "Test solde insuffisant"
}
\end{lstlisting}

\begin{dangerbox}
\textbf{Résultat attendu :} HTTP 400\\
Corps : \texttt{\{"detail": "Solde insuffisant pour effectuer ce retrait."\}}
\end{dangerbox}

\subsubsection{Étape 7 — Consulter l'historique (CT-22)}

\begin{lstlisting}[style=bash, caption={Appel GET /accounts/\{id\}/transactions}]
GET http://localhost:8000/accounts/a3f1b2c4-.../transactions
\end{lstlisting}

\begin{successbox}
\textbf{Résultat attendu :} HTTP 200 avec 3 transactions dans l'ordre chronologique :\\
1. Dépôt initial (10 000)\\
2. Dépôt agence (5 000) → balance 15 000\\
3. Retrait DAB (2 000) → balance 13 000
\end{successbox}

\subsubsection{Étape 8 — Tester un ID inexistant (CT-09)}

\begin{lstlisting}[style=bash, caption={Test CT-09 : compte inexistant}]
GET http://localhost:8000/accounts/00000000-0000-0000-0000-000000000000
\end{lstlisting}

\begin{dangerbox}
\textbf{Résultat attendu :} HTTP 404\\
Corps : \texttt{\{"detail": "Compte introuvable."\}}
\end{dangerbox}

% =========================================================
\section{Documentation Swagger}
% =========================================================

\subsection{Accès à l'interface}

L'interface Swagger UI est accessible à l'adresse \texttt{http://localhost:8000/docs} une fois le serveur lancé. Elle permet de tester toutes les fonctions de l'API directement depuis le navigateur, sans outil externe.

\subsection{Exemples interactifs dans Swagger}

L'API intègre des \textbf{exemples nommés} dans Swagger pour chaque endpoint de création et de transaction. Ces exemples apparaissent dans un menu déroulant \textit{«~Examples~»} permettant de basculer entre les cas de test.

\begin{center}
\begin{tabular}{lll}
\toprule
\textbf{Endpoint} & \textbf{Exemples disponibles} \\
\midrule
\texttt{POST /accounts} & CT-01, CT-02, CT-03, CT-04, CT-05 \\
\texttt{POST /\{id\}/deposit} & CT-10, CT-11, CT-12, CT-13 \\
\texttt{POST /\{id\}/withdraw} & CT-15, CT-16, CT-17, CT-18, CT-19 \\
\bottomrule
\end{tabular}
\end{center}

\subsection{Comment utiliser Swagger pour le test manuel}

\begin{enumerate}
  \item Lancer l'API : \texttt{uvicorn app.main:app -{}-reload}
  \item Ouvrir \texttt{http://localhost:8000/docs} dans le navigateur
  \item Cliquer sur un endpoint (ex : \textit{POST /accounts})
  \item Cliquer sur le bouton \textbf{«~Try it out~»}
  \item Sélectionner un exemple dans le menu déroulant \textbf{«~Examples~»}
  \item Cliquer sur \textbf{«~Execute~»}
  \item Observer le code HTTP et le corps de la réponse dans la section \textit{Responses}
\end{enumerate}

% =========================================================
\section{Lancement et déploiement}
% =========================================================

\subsection{Lancement en local}

\begin{lstlisting}[style=bash, caption={Lancement de l'API en local}]
# 1. Installer les dependances
pip install -r requirements.txt

# 2. Lancer le serveur (mode developpement avec rechargement automatique)
uvicorn app.main:app --reload

# 3. L'API est disponible sur :
#    http://localhost:8000/
#    http://localhost:8000/docs  (Swagger UI)
\end{lstlisting}

\subsection{Contenu de requirements.txt}

\begin{lstlisting}[style=bash, caption={requirements.txt}]
fastapi
uvicorn[standard]
\end{lstlisting}

\subsection{Déploiement sur Render}

Le fichier \texttt{render.yaml} à la racine du projet configure le déploiement automatique sur la plateforme cloud \textbf{Render.com}. À chaque push sur la branche \texttt{main} de GitHub, le service est automatiquement mis à jour.

% =========================================================
\section{Conclusion}
% =========================================================

Ce projet implémente une \textbf{API REST complète} pour la gestion de comptes bancaires, respectant toutes les exigences fonctionnelles et non fonctionnelles définies dans le cahier des charges.

\subsection*{Points clés de l'implémentation}

\begin{itemize}[noitemsep]
  \item \textbf{Séparation des responsabilités} : 3 modules distincts avec des rôles clairs
  \item \textbf{Validation automatique} : Pydantic assure la validité de toutes les données entrantes
  \item \textbf{Persistance fiable} : fichier JSON protégé par un verrou thread-safe
  \item \textbf{Documentation interactive} : 23 cas de test directement exécutables dans Swagger
  \item \textbf{Codes HTTP corrects} : 201 (création), 200 (succès), 400 (logique métier), 404 (non trouvé), 422 (validation)
\end{itemize}

\subsection*{Ce que ce projet ne fait pas encore (limites)}

\begin{itemize}[noitemsep]
  \item Pas d'authentification (JWT ou sessions)
  \item Pas de base de données relationnelle (SQLite, PostgreSQL)
  \item Pas de virement entre comptes
  \item Pas de séparation des rôles client/administrateur
\end{itemize}

Ces éléments constituent des évolutions naturelles pour une version future du système.

\vfill
\begin{center}
\rule{0.5\linewidth}{0.5pt}\\[0.3cm]
{\small KOA ESSAMA PAULIN BRICE — Matricule 25I2255 — Licence 3 ICT\\
Université de Yaoundé I — Avril 2026}
\end{center}

\end{document}
